\section{서론}

\subsection*{배경과 문제}
제조·물류 현장에 로봇팔, AMR, 시각-언어-행동(VLA) 모델~\cite{rt2} 기반 Physical AI의 도입이 늘고 있다. 이들은 실제 공간에서
행동하므로 소프트웨어 정확도만으로는 품질을 보장할 수 없고, 물체 위치 변화나 장애물 추가,
작업자 접근에 따라 현장에서 실패할 수 있다. 현재의 검증은 사람이 예상 가능한 케이스를
수작업으로 정의하거나 실제 장비에서 제한적으로 반복하는 방식에 의존한다. 이 방식은
테스트 작성에 시간이 들고, 충돌·안전 위험이 있으며, 드문 실패나 PASS/FAIL 경계 조건을
체계적으로 찾기 어렵다. 현장 배치가 바뀔 때마다 테스트를 다시 작성해야 하는 것도 문제다.

\subsection*{목표 지표}
과제의 평가 지표와 목표는 표~\ref{tab:metric}과 같다. 세 지표 모두 같은 시험 예산에서 Random 탐색과 비교한다.
\begin{table*}[!t]
\centering\small
\caption{과제 평가 지표와 목표}\label{tab:metric}
\begin{tabularx}{\textwidth}{@{}p{38mm}Yp{40mm}@{}}
\toprule
지표 & 정의 & 목표 \\
\midrule
실패 발견율 (FDR) & 유효 테스트 중 작업 실패를 발견한 비율, $N_{\mathrm{FAIL}}/N_{\mathrm{VALID}}$ & 30\% 이상 \\
실패 탐색 성능 향상 (Gain) & 동일 테스트 예산에서 Random 대비 실패 발견 향상 & 20\% 이상 \\
실패 유형 다양성 & 사전 정의한 주요 실패 유형 6종 중 자동 발견한 유형 수 & 6종 중 4종 이상 \\
\bottomrule
\end{tabularx}
\end{table*}

\subsection*{기여}
\begin{enumerate}
\item 장면 파라미터로부터 3D 장면 그래프, 이동 지도, 물리 시뮬레이션 모델을 일관되게 생성하고 실행 전에 유효성을 검사하는 환경 생성기.
\item 이전 실행의 행동, 로봇 기능 호출 결과, 측정값을 근거로 다음 테스트 환경을 제안하는 LLM 기반 AFS와 후보 검증기.
\item 실행 유효성, 목표 달성, 실패 유형을 단계적으로 판정하는 평가기와 6종/12종 실패 분류 체계.
\item 로봇팔(VLA)과 휴머노이드(LLM Agent) 두 플랫폼에서 Random 대비 AFS의 실패 탐색 효율을 검증.
\item 발견한 실패·성공 사례를 다시 실행하여 이전 결과와 비교하는 회귀 테스트 실행기.
\end{enumerate}
